\chapter{Parcours professionnel détaillé}

Mon parcours professionnel de 25 ans révèle une progression cohérente de l'expertise technique industrielle vers l'enseignement supérieur. Cette trajectoire s'articule autour de quatre phases distinctes qui démontrent une montée en compétences constante et une adaptation permanente aux évolutions technologiques et pédagogiques.

\section{Vue d'ensemble chronologique}

\begin{center}
\textit{Synthèse : 3 ans d'autoformation, 2 ans de relation client, 1 an de développement, 12 ans de management technique chez un leader du service numérique, et 5 ans d'enseignement en lycée.}
\end{center}

\section{Team Leader chez TOTAL/Stéria puis SopraStéria (2007-2019)}

\subsection{Contexte et développement professionnel}

Cette expérience de 12 années constitue la période la plus formatrice de mon parcours professionnel. L'évolution de TOTAL/Stéria vers SopraStéria, leader européen des services numériques (46 000 collaborateurs), m'a permis d'évoluer dans un environnement international avec des projets d'envergure pour grands comptes.

\textbf{Ma fonction :} Team Leader technique, avec encadrement d'équipes de développeurs, gestion de projets clients, formation technique des collaborateurs, et interface client-équipe technique.

\textbf{Évolution progressive :} De développeur senior vers team leader puis responsable technique d'équipe. J'ai managé des projets pour des clients majeurs (banques, assurances, industrie) dans un contexte international. Cette période a été marquée par la formation de nombreux collaborateurs, la mise en place de méthodologies Agile et la gestion directe de la relation client.

\subsection{Évolution des responsabilités (2007-2019)}

Durant cette période (2007-2019), mes responsabilités ont évolué de l'expertise technique individuelle vers le management opérationnel puis stratégique. En tant que Lead Technique, j'ai assuré l'architecture des solutions et l'innovation, tout en continuant à mentorer les équipes juniors.

\subsection{Synthèse des compétences développées}

Cette expérience m'a permis de développer une quadruple compétence :
\begin{itemize}[leftmargin=*]
    \item \textbf{Management} : Encadrement d'équipes, recrutement, évaluation et développement des talents (leadership transformationnel).
    \item \textbf{Gestion de Projet} : Pilotage de projets complexes, planification stratégique et maîtrise des méthodologies Agile/Scrum.
    \item \textbf{Relation Client} : Interface technique, conseil stratégique, communication et gestion de crise.
    \item \textbf{Expertise Technique} : Architecture logicielle, veille technologique, revue de code et garantie de la qualité technique.
\end{itemize}

Ces 12 années d'expérience managériale continue, avec un fort taux de réussite des projets et une évolution significative des collaborateurs formés, constituent le cœur de mon expertise professionnelle.



Cette expérience correspond directement aux compétences de leadership et de pratique réflexive du Master MEEF par le développement de compétences de leadership, formation d'équipes et pratique réflexive \parencite{referentiel_meef}.

\section{Enseignant NSI/SNT (Éducation Nationale) 2020-2025}

\subsection{Développement professionnel}

\textbf{Contexte :} Reconversion vers l'enseignement après obtention CAPES NSI et CAPET SI en 2020, affectation lycée général et technologique, classes de Seconde (SNT) et Première/Terminale (NSI).

\textbf{Ma fonction :} Professeur certifié, enseignement Numérique et Sciences Informatiques et Sciences Numériques et Technologie, 18h/semaine, suivi de 120 élèves/an, participation vie établissement.

\textbf{Transition réussie :} Adaptation des compétences managériales vers la pédagogie. Développement de séquences pédagogiques innovantes, utilisation d'outils numériques variés, différenciation pédagogique. Participation active aux formations académiques et aux groupes de travail NSI.



\subsection{Synthèse des compétences enseignantes}

Mon activité d'enseignant s'articule autour de trois pôles :
\begin{itemize}[leftmargin=*]
    \item \textbf{Pédagogie} : Enseignement SNT/NSI, différenciation pédagogique, évaluation formative et gestion de classe.
    \item \textbf{Numérique} : Intégration des outils numériques (ENT, Moodle), programmation (Python) et innovation pédagogique.
    \item \textbf{Formation Continue} : Participation active aux formations académiques, groupes de travail et veille pédagogique constante.
\end{itemize}

\section{Enseignant PRCE (2025-Présent) et Directeur Adjoint au Numérique (Septembre 2026)}

\subsection{Fonctions et gouvernance institutionnelle}

\textbf{Contexte :} Affectation en tant qu'enseignant PRCE au département Techniques de Commercialisation (TC) de l'IUT du Havre depuis septembre 2025, et prise de fonction en tant que \textbf{Directeur Adjoint dédié aux questions du numérique} pour l'ensemble de l'IUT à compter de \textbf{septembre 2026}.

\textbf{Responsabilités clés :}
\begin{itemize}[leftmargin=*]
    \item \textbf{Pilotage stratégique du numérique} : Définition des orientations numériques de l'établissement, gouvernance des équipements informatiques et accompagnement de la transformation pédagogique numérique.
    \item \textbf{Intégration de l'IA générative} : Cadrage éthique, pédagogique et technique de l'intégration des assistants IA pour la formation et l'évaluation au sein des départements de l'IUT.
    \item \textbf{Ingénierie de formation et enseignements} : Responsable des modules de \enquote{Ressources et Culture Numérique} (BUT TC) et soutien aux équipes pédagogiques.
\end{itemize}